```latex
\documentclass{article}
\usepackage{amsmath,amssymb,amsthm}

\newtheorem{lemma}{Lemma}
\newtheorem{theorem}{Theorem}

\theoremstyle{remark}
\newtheorem{remark}{Remark}

\begin{document}

\title{A cubic identity in hyperfields}
\author{}
\date{}

\maketitle

Let $\mathbb{H}$ be a hyperfield with multivalued addition $\boxplus$. 
We use the standard hyperfield axioms: $\boxplus$ is commutative and associative, 
$0\boxplus x=\{x\}$ for all $x$, every $x$ has a unique additive inverse $-x$ with 
$0\in x\boxplus(-x)$, and multiplication distributes over $\boxplus$.

\begin{lemma}
\label{lem:zero-preservation}
Let $A_1,\dots,A_n\subseteq \mathbb{H}$ be nonempty subsets such that $0\in A_i$ 
for every $i$. Suppose that a formal hyperaddition
\[
A_1\boxplus \cdots \boxplus A_n
\]
is simplified by deleting a pair of opposite terms $u$ and $-u$. 
If $u\neq 0$, delete only those two terms; if $u=0$, replace the two zero terms 
by one zero term. Then the resulting hyperaddition still contains $0$.
\end{lemma}

\begin{proof}
It suffices to prove the claim for one deletion. Let $A_i'$ be the subsets 
obtained after the deletion.

If $u\neq 0$, then neither $u$ nor $-u$ is $0$, so every $A_i'$ still contains 
$0$. Hence
\[
0\in \underbrace{0\boxplus\cdots\boxplus 0}_{n\text{ terms}}
\subseteq A_1'\boxplus\cdots\boxplus A_n'.
\]
If $u=0$, then the two deleted terms are both $0$, and we replace them by one 
zero term. Again every $A_i'$ contains $0$, so the same argument applies.
\end{proof}

\begin{theorem}
\label{thm:cubic}
Let $\mathbb{H}$ be a hyperfield and let $x,y,z\in \mathbb{H}$. If
\[
0\in x\boxplus y\boxplus z,
\]
then
\[
0\in x^3\boxplus y^3\boxplus z^3\boxplus (-xyz)\boxplus (-xyz)\boxplus (-xyz).
\]
\end{theorem}

\begin{proof}
Assume that $0\in x\boxplus y\boxplus z$. Define six subsets
\[
\begin{aligned}
A_1 &= x^2(x\boxplus y\boxplus z),\\
A_2 &= y^2(x\boxplus y\boxplus z),\\
A_3 &= z^2(x\boxplus y\boxplus z),\\
A_4 &= (-xy)(x\boxplus y\boxplus z),\\
A_5 &= (-yz)(x\boxplus y\boxplus z),\\
A_6 &= (-zx)(x\boxplus y\boxplus z).
\end{aligned}
\]
Since $0\in x\boxplus y\boxplus z$, multiplication by any element gives
\[
0\in A_i \qquad (i=1,\dots,6).
\]
By distributivity and associativity of $\boxplus$, the full expansion
\[
A_1\boxplus A_2\boxplus A_3\boxplus A_4\boxplus A_5\boxplus A_6
\]
is exactly the expansion of the product
\[
(x\boxplus y\boxplus z)\left(x^2\boxplus y^2\boxplus z^2\boxplus (-xy)\boxplus (-yz)\boxplus (-zx)\right).
\]
Explicitly,
\[
\begin{aligned}
A_1 &= x^3\boxplus x^2y\boxplus x^2z,\\
A_2 &= xy^2\boxplus y^3\boxplus y^2z,\\
A_3 &= xz^2\boxplus yz^2\boxplus z^3,\\
A_4 &= -x^2y\boxplus -xy^2\boxplus -xyz,\\
A_5 &= -xyz\boxplus -y^2z\boxplus -yz^2,\\
A_6 &= -x^2z\boxplus -xyz\boxplus -xz^2.
\end{aligned}
\]
Now delete the following opposite pairs:
\[
\begin{aligned}
&x^2y \text{ and } -x^2y,\\
&xy^2 \text{ and } -xy^2,\\
&x^2z \text{ and } -x^2z,\\
&xz^2 \text{ and } -xz^2,\\
&y^2z \text{ and } -y^2z,\\
&yz^2 \text{ and } -yz^2.
\end{aligned}
\]
After these deletions, the remaining terms are precisely
\[
x^3,\quad y^3,\quad z^3,\quad -xyz,\quad -xyz,\quad -xyz.
\]
Since each $A_i$ contains $0$, Lemma~\ref{lem:zero-preservation} applies after each 
deletion. Therefore the final hyperaddition contains $0$, i.e.
\[
0\in x^3\boxplus y^3\boxplus z^3\boxplus (-xyz)\boxplus (-xyz)\boxplus (-xyz).
\]
\end{proof}

\begin{remark}
The argument is the hyperfield analogue of the ordinary polynomial identity
\[
x^3+y^3+z^3-3xyz=(x+y+z)(x^2+y^2+z^2-xy-yz-zx).
\]
The lemma is needed because in a hyperfield an opposite pair $u\boxplus(-u)$ need not 
collapse to $\{0\}$; however, the presence of zero in each factor of the expansion 
is enough to preserve the desired zero after the usual cancellations.
\end{remark}

\end{document}
```